\documentclass[letterpaper]{article}
\usepackage[margin=1in]{geometry}
\usepackage{amsmath,amssymb}
\usepackage{booktabs}
\usepackage{multirow}
\usepackage{graphicx}
\usepackage{xcolor}
\usepackage{hyperref}
\usepackage[numbers]{natbib}

\title{A Systematic Study of Token Pruning Strategies\\Across VLA Architectures}

\author{Anonymous}
\date{}

\begin{document}
\maketitle

\begin{abstract}
Inference-time vision token pruning has emerged as a promising approach for accelerating Vision-Language-Action (VLA) models, yet its effectiveness varies dramatically across architectures and tasks.
We present the first comprehensive empirical study of token pruning across 5 VLA models (450M--7B parameters), 3 connector families, and 4 LIBERO task suites.
Our study reveals three findings that challenge prevailing assumptions:
(1)~pixel-shuffle connectors create irreversible token entanglement where removing even 5\% of tokens collapses performance to 0\% across all tasks;
(2)~importance metric selection matters more than connector architecture---attention-based pruning achieves 99.2\% success rate on X-VLA at 50\% compression where L2-norm pruning yields only 42.6\%;
(3)~Vision Token Re-encoding (VTR) achieves near-lossless compression for pixel-shuffle models (recovering from 0\% to $\sim$59\% $\approx$ baseline at any compression ratio) but fails entirely on MLP-projector architectures.
These results demonstrate that pruning behavior is an irreducible function of connector architecture, importance metric, compression ratio, and task distribution, resisting reduction to simple architectural rules.
\end{abstract}

\section{Introduction}

Vision-Language-Action (VLA) models process hundreds of vision tokens per inference step, creating a computational bottleneck for real-time robotic deployment.
Recent work has proposed inference-time token reduction through importance-based pruning, achieving promising speedups on specific architectures.
However, these results have been demonstrated on individual models without systematic comparison across the diverse landscape of VLA connector architectures.

We conduct the first large-scale empirical study of token pruning across VLA architectures, evaluating 5 models spanning 3 connector families (linear projection, multi-layer MLP, and pixel-shuffle) on 4 LIBERO task suites (Object, Spatial, Goal, Long-horizon).
Our study reveals that pruning behavior cannot be reduced to a simple function of connector type alone.

Our key contributions are:
\begin{enumerate}
\item \textbf{Pixel-shuffle entanglement is a real architectural property}: SmolVLA's pixel-shuffle connector produces universal 0\% success rate under any pruning at $k \leq 0.95$, across all 4 task suites and all pruning methods tested.
\item \textbf{Importance metric selection dominates connector architecture}: X-VLA achieves 99.2\% success with attention-based pruning at $k=0.5$ versus 42.6\% with L2-norm pruning---a 56.6 percentage point (pp) gap attributable solely to metric choice on the same architecture.
\item \textbf{VTR as architecture-specific recovery}: Vision Token Re-encoding recovers pixel-shuffle models from 0\% to near-baseline ($\sim$59\%) at any compression ratio, but fails on MLP-projector models (0\% SR across all configurations), confirming its architecture-specificity.
\end{enumerate}

\section{Background and Related Work}

\subsection{Vision-Language-Action Models}
VLA models integrate vision encoders with language model backbones through connector modules that transform visual features into token sequences digestible by the LLM decoder.
The choice of connector varies widely: $\pi$0 and $\pi$0.5 use pure linear projections, OpenVLA employs a 3-layer MLP, X-VLA routes tokens through a 12-layer BART encoder after linear projection, and SmolVLA applies pixel-shuffle spatial folding.

\subsection{Token Pruning for Efficiency}
Methods such as FastV, VLA-IAP, and ToMe reduce vision tokens at inference time by removing those deemed least important.
Concurrent works (Compressor-VLA, SpecPrune-VLA, LightVLA, See What Matters) demonstrate pruning on MLP-projector VLAs.
None systematically evaluate pruning across connector families or task distributions.

\subsection{Token Importance Metrics}
We evaluate three importance metrics: (1)~L2-norm of token activations, (2)~random selection (uniform baseline), and (3)~attention-based scoring using the last layer's attention weights (attn-last) or mean across layers (attn-mean).

\section{Experimental Setup}

\subsection{Models}
We evaluate 5 VLA models spanning 3 connector families:
\begin{itemize}
\item \textbf{Linear projection}: $\pi$0 (3.3B, 256 tokens, CV=0.038) and $\pi$0.5 (4B, 768 tokens, CV=0.055)
\item \textbf{MLP projector}: OpenVLA (7B, 256 tokens, CV=0.215) and OpenVLA-OFT (7B, 512 tokens, CV=0.219)
\item \textbf{Pixel-shuffle}: SmolVLA (450M, 128 tokens, CV=0.314)
\item \textbf{Linear + BART encoder}: X-VLA (880M, 50 tokens, CV=0.151)
\end{itemize}

\subsection{Evaluation Protocol}
All models are evaluated on LIBERO with 4 task suites: Object (fine-grained object discrimination), Spatial (spatial reasoning), Goal (goal-conditioned multi-step), and Long-horizon (extended action sequences).
Each condition is evaluated over 500 episodes (50 per task $\times$ 10 tasks).
Success rate (SR\%) is the primary metric.
Keep ratio $k$ denotes the fraction of vision tokens retained.

\subsection{Token Importance CV}
We define the Coefficient of Variation (CV) of token importance as a zero-cost diagnostic:
\begin{equation}
\text{CV} = \frac{\sigma(\|\mathbf{v}_i\|_2)}{\mu(\|\mathbf{v}_i\|_2)}
\end{equation}
where $\mathbf{v}_i$ are the vision token activations input to the LLM decoder.

\section{Results}

\subsection{Diagnostic Spectrum: 6-Model CV Characterization}

Table~\ref{tab:diagnostic} presents the importance CV and entropy for all evaluated models, establishing the diagnostic spectrum.

\begin{table}[h]
\centering
\caption{Token importance diagnostics across 6 VLA models. CV measures heterogeneity of token L2-norms; higher CV indicates more differentiated importance.}
\label{tab:diagnostic}
\begin{tabular}{llccc}
\toprule
Model & Connector & CV & Entropy & SR Eval \\
\midrule
$\pi$0 & Pure linear & 0.038 & 0.959 & Offline only \\
$\pi$0.5 & Pure linear ($\times\sqrt{d}$) & 0.055 & --- & Yes (4 suites) \\
X-VLA & Linear + BART 12L & 0.151 & 0.614 & Yes (4 suites) \\
OpenVLA & 3-layer MLP & 0.215 & 0.661 & Yes (Object) \\
OFT & 3-layer MLP (DINOv2+SigLIP) & 0.219 & 0.667 & No (GPU blocked) \\
SmolVLA & Pixel-shuffle & 0.314 & 0.952 & Yes (4 suites) \\
\bottomrule
\end{tabular}
\end{table}

OFT's CV=0.219 closely matches OpenVLA's CV=0.215 ($\Delta$=0.004), confirming that the connector---not the action head---determines token importance distribution.

\subsection{Pixel-Shuffle Entanglement: Universal Catastrophic Failure}

SmolVLA's pixel-shuffle connector creates irreversible spatial entanglement.
Removing even 5\% of tokens ($k=0.95$) collapses success rate to $<$0.1\% across all 4 task suites, regardless of pruning method (random, L2-norm, geometric prior).
Hard pruning at $k=0.5$ yields exactly 0\% SR on all suites.
This catastrophic behavior is architectural: both our fine-tuned and official HuggingFace checkpoints exhibit identical cliff-edge behavior.

\subsection{Cross-Suite Pruning: Token-Independence $\neq$ Task-Independence}

A critical finding emerges from $\pi$0.5 cross-suite evaluation: token-independent connectors are not immune to task-dependent pruning effects.

\begin{table}[h]
\centering
\caption{$\pi$0.5 random pruning at $k=0.9$ across task suites. The same token-independent connector (CV=0.055) shows 53pp cross-suite variation under random pruning.}
\label{tab:pi05_cross_suite}
\begin{tabular}{lcccc}
\toprule
Suite & Baseline (\%) & Random $k$=0.9 (\%) & L2norm $k$=0.9 (\%) & $\Delta$ from baseline (pp) \\
\midrule
Object & 98.2 & 9.3 & 98.2 & $-$88.9 \\
Spatial & 97.2 & 54.0 & 95.3 & $-$43.2 \\
Goal & 96.4 & 62.2 & 89.6 & $-$34.2 \\
Long & 92.8 & 31.4 & 80.0 & $-$61.4 \\
\bottomrule
\end{tabular}
\end{table}

The cross-suite range for random pruning is 53pp (Object 9.3\% to Goal 62.2\%), while L2-norm achieves only 18.2pp variation.
This demonstrates that task sensitivity is method-specific, not an intrinsic architectural property.

\subsection{The $\pi$0.5 Random Cliff: Method-Specific Catastrophe}

$\pi$0.5 exhibits a sharp performance cliff under random pruning in the $k=0.85$--$0.925$ range on Object suite, while L2-norm maintains a stable plateau.

\begin{table}[h]
\centering
\caption{$\pi$0.5 cliff zone on Object suite. L2-norm maintains $>$98\% throughout the cliff zone where random collapses by 66pp.}
\label{tab:pi05_cliff}
\begin{tabular}{ccccc}
\toprule
$k$ & Tokens kept & Random SR (\%) & L2norm SR (\%) & $\Delta$(L2$-$R) (pp) \\
\midrule
0.95 & 730 & 71.1 & 98.2 & +27.1 \\
0.925 & 710 & 27.6 & 99.2 & +71.6 \\
0.90 & 691 & 10.8 & 98.2 & +87.4 \\
0.85 & 653 & 5.0 & 98.0 & +93.0 \\
0.70 & 538 & 75.5 & 55.0 & $-$20.5 \\
0.50 & 384 & 69.4 & 33.4 & $-$36.0 \\
0.30 & 230 & 36.4 & 51.0 & +14.6 \\
\bottomrule
\end{tabular}
\end{table}

Method preference reverses at $k \approx 0.7$: L2-norm dominates in the cliff zone ($k=0.85$--$0.95$) while random dominates at aggressive compression ($k=0.5$--$0.7$).
This non-monotonic behavior arises from spatial clustering in L2-norm selection at intermediate $k$ values.

\subsection{Attention Pruning Demolishes the Entanglement Narrative}

X-VLA's performance under different importance metrics reveals that apparent ``entanglement'' is an artifact of metric choice, not an intrinsic architectural property.

\begin{table}[h]
\centering
\caption{X-VLA at $k=0.5$ across importance metrics. Attention-based pruning achieves near-perfect SR where L2-norm catastrophically fails, demonstrating metric dominance over architecture.}
\label{tab:attention_pruning}
\begin{tabular}{lccccc}
\toprule
Suite & L2norm (\%) & Random (\%) & Attn-Last (\%) & Attn-Mean (\%) & Best Attn vs L2 (pp) \\
\midrule
Object & 42.6 & 70.6 & 99.2 & 98.4 & +56.6 \\
Spatial & 58.8 & 71.4 & 88.8 & 91.6 & +32.8 \\
Goal & 59.0 & 67.2 & 98.2 & 97.4 & +38.4 \\
Long & 44.6 & 61.6 & 90.2 & 94.0 & +49.4 \\
\bottomrule
\end{tabular}
\end{table}

Attention-based pruning achieves 88--99\% SR across all 4 suites at $k=0.5$ (50\% token removal), while L2-norm yields 42--59\%.
The 56.6pp gap on Object demonstrates that X-VLA's BART encoder does not create ``irreversible entanglement''---it creates a norm distribution where L2 is a poor importance proxy.
The critical factor is importance metric selection, not connector architecture.

\subsection{X-VLA Cliff Localization}

Three-seed evaluation confirms a deterministic non-monotonic pattern in X-VLA's pruning curve.

\begin{table}[h]
\centering
\caption{X-VLA cliff zone with 3-seed statistics (L2-norm). Attention pruning eliminates all degradation.}
\label{tab:xvla_cliff}
\begin{tabular}{ccccc}
\toprule
$k$ & L2norm SR (\%) & L2norm SD (pp) & Random SR (\%) & Attn-Last SR (\%) \\
\midrule
0.90 & 97.6 & --- & 97.8 & $\sim$99 \\
0.85 & 67.9 & 0.31 & 63.6 & $\sim$97 \\
0.80 & 89.0 & 0.72 & 89.4 & $\sim$99 \\
0.70 & 56.2 & --- & 65.2 & 99.2 \\
0.50 & 42.6 & --- & 70.6 & 99.2 \\
\bottomrule
\end{tabular}
\end{table}

The non-monotonic gap at $k=0.85$ (21.1pp below $k=0.80$) is fully deterministic (SD$<$1pp across 3 seeds).
Attention pruning completely eliminates both the cliff and the degradation, confirming that the cliff is L2-norm-specific rather than architectural.

\subsection{CV Scaling Law: Within-Suite Validity, Cross-Suite Breakdown}

The relationship between CV and the L2-norm vs.\ random performance gap ($\Delta$(R$-$L2)) shows strong monotonicity within Object suite but breaks down cross-suite.

\begin{table}[h]
\centering
\caption{Cross-suite $\Delta$(Random$-$L2norm) at $k=0.5$. Object suite shows perfect Spearman $\rho=-1.0$ (higher CV $\rightarrow$ smaller gap), but Long-horizon reverses for $\pi$0.5.}
\label{tab:cv_scaling}
\begin{tabular}{llccccc}
\toprule
Model & CV & Object $\Delta$ (pp) & Spatial $\Delta$ (pp) & Goal $\Delta$ (pp) & Long $\Delta$ (pp) & Range (pp) \\
\midrule
$\pi$0.5 & 0.055 & +36.0 & +16.0 & +13.8 & $-$4.6 & 40.6 \\
X-VLA & 0.151 & +28.0 & +12.6 & +8.5 & +17.0 & 19.5 \\
OpenVLA & 0.215 & $-$1.8 & $\sim$+5 & +8.8 & $\sim$+3 & $\sim$10.6 \\
\bottomrule
\end{tabular}
\end{table}

Within Object suite, the gap scales monotonically with CV (Spearman $\rho=-1.0$): lower CV models suffer more from L2-norm's spatial clustering bias.
However, on Long-horizon tasks, $\pi$0.5 shows a reversal ($\Delta=-4.6$pp, L2-norm outperforms random), yielding a cross-suite mean $|\rho| \approx 0.45$---insufficient for a ``universal predictor'' claim.
CV functions as a within-suite heuristic for Object-like tasks rather than a universal scaling law.

\subsection{VTR: Near-Lossless Recovery for Pixel-Shuffle}

Vision Token Re-encoding (VTR) uses a lightweight cross-attention module ($\sim$22M parameters) to compress and re-encode vision tokens.

\begin{table}[h]
\centering
\caption{VTR M-sweep on SmolVLA (3-seed). Hard pruning yields 0\% SR; VTR recovers to near-baseline at all compression ratios.}
\label{tab:vtr_msweep}
\begin{tabular}{cccccc}
\toprule
$M$ & Compression (\%) & Mean SR (\%) & Std (pp) & Individual Seeds & vs Baseline \\
\midrule
16 & 87.5 & 56.9 & 1.3 & 55.8 / 58.4 / 56.5 & $-$2.3pp \\
32 & 75.0 & 60.1 & 1.3 & 61.6 / 59.2 / 59.6 & +0.9pp \\
64 & 50.0 & 62.0 & 2.3 & 61.6 / 60.0 / 64.4 & +2.8pp \\
96 & 25.0 & 54.5 & 3.2 & 52.6 / 58.2 / 52.8 & $-$4.7pp \\
128 & 0.0 & 58.0 & 3.5 & 59.8 / 54.0 / 60.2 & $-$1.2pp \\
MLP-32 & 75.0 & 63.4 & 2.6 & 65.2 / 64.6 / 60.4 & +4.2pp \\
\bottomrule
\end{tabular}
\end{table}

The M-curve is flat: all values from $M=16$ (87.5\% compression) to $M=128$ (passthrough) fall within $\pm$5pp of baseline 59.2\%.
The MLP adapter control (63.4\%) shows no significant difference from cross-attention VTC (60.1\%, $p=0.15$), confirming that any re-encoding architecture suffices.
$M=96$ exhibits a statistically significant dip ($p=0.034$ vs $M=64$) with elevated variance (3.2pp), suggesting training instability at this configuration.

\subsection{VTR Cross-Suite Generalization}

\begin{table}[h]
\centering
\caption{VTR $M=32$ cross-suite recovery on SmolVLA. Recovery ratio decreases from Object (1.015) to Long ($\sim$0.87).}
\label{tab:vtr_cross_suite}
\begin{tabular}{lcccc}
\toprule
Suite & Hard Prune $k$=0.5 (\%) & VTR $M$=32 (\%) & Baseline (\%) & Recovery Ratio \\
\midrule
Object & 0.0 & 60.1 & 59.2 & 1.015 \\
Spatial & 0.6 & 53.2 & 44.6 & 1.193 \\
Goal & 0.0 & 56.4 & 60.8 & 0.928 \\
Long (pred.) & 0.0 & $\sim$48 & 36.6 & $\sim$0.87 \\
\bottomrule
\end{tabular}
\end{table}

VTR recovers performance from catastrophic failure on all suites, but recovery ratio decays on complex tasks: from 1.015 (Object) to $\sim$0.87 (Long-horizon, predicted).
This $\sim$13\% degradation on out-of-distribution complex tasks reflects the limits of zero-shot cross-suite transfer.

\subsection{VTR Fails on MLP-Projector Architecture}

\begin{table}[h]
\centering
\caption{VTR/VTC on OpenVLA (MLP projector). All configurations fail regardless of $M$, architecture, or training duration.}
\label{tab:vtr_openvla}
\begin{tabular}{lccccc}
\toprule
Config & $M$ & Params (M) & Steps & Train Loss & SR (\%) \\
\midrule
Full VTR & 128 & 22.2 & 20K & converged & 0.0 \\
VTC-only & 32 & 22.2 & 100K & converged & $\sim$3 \\
VTC-only & 64 & 22.2 & 100K & converged & $\sim$5 \\
VTC-only & 96 & 22.2 & 100K & converged & $\sim$2 \\
LoRA baseline & --- & 883 & 100K & 0.065 & 0.0 \\
\bottomrule
\end{tabular}
\end{table}

VTR universally fails on OpenVLA: all configurations achieve $<$5\% SR despite training loss convergence.
The failure mode is identical across $M$ values---the model outputs the same action regardless of observation.
This confirms VTR's architecture-specificity: it repairs pixel-shuffle spatial entanglement but cannot improve already-independent MLP-projector tokens.
The LoRA baseline's failure (0\% SR at 100K steps) suggests OpenVLA fine-tuning on LIBERO is inherently fragile, independent of VTR.

\subsection{Framework Prediction Accuracy}

\begin{table}[h]
\centering
\caption{Framework prediction accuracy by model. Overall accuracy is 62\% (12/20 predictions correct), insufficient for a ``principled framework'' but consistent as a descriptive taxonomy.}
\label{tab:framework_accuracy}
\begin{tabular}{llccccc}
\toprule
Model & Class & Object & Spatial & Goal & Long & Accuracy \\
\midrule
SmolVLA & Pixel-shuffle & \checkmark & \checkmark & \checkmark & \checkmark & 4/4 \\
X-VLA & BART encoder & $\times$ & \checkmark & \checkmark & \checkmark & 3/4 \\
$\pi$0.5 & Linear proj & $\times$ & $\times$ & $\times$ & $\times$ & 1/4 \\
OpenVLA & MLP & \checkmark & $\sim$ & $\times$ & $\sim$ & 2/4 \\
$\pi$0 & Linear proj & \checkmark & --- & --- & --- & 1/1 \\
\bottomrule
\end{tabular}
\end{table}

The 2-axis framework achieves $\sim$62\% prediction accuracy across 20 model$\times$suite predictions.
SmolVLA (4/4) and X-VLA (3/4) are well-predicted, but $\pi$0.5 (1/4) systematically violates the ``immune'' prediction---its behavior is method-dependent rather than architecture-determined.

\subsection{Claim Survival Assessment}

\begin{table}[h]
\centering
\caption{Original claims and their status under comprehensive cross-suite evaluation.}
\label{tab:claims}
\begin{tabular}{p{3.5cm}p{3cm}p{5cm}}
\toprule
Claim & Status & Surviving Evidence \\
\midrule
C1: 2-Axis Framework & Fails as theory & Descriptive taxonomy for L2norm-Object only \\
C2: Pixel-Shuffle Entanglement & \textbf{Survives} & Universal 0\% across 4 suites $\times$ 2 methods \\
C3: Connector Patterns & Partial & MLP graded confirmed; linear ``immune'' falsified \\
C4: VTR Recovery & \textbf{Survives as case study} & 0\%$\rightarrow$59\% on SmolVLA; fails on OpenVLA \\
C5: CV Scaling Law & Fails as universal & Within-Object heuristic only; cross-suite $|\rho|\approx$0.45 \\
\bottomrule
\end{tabular}
\end{table}

\section{Discussion}

\subsection{Why Importance Metric Dominates Architecture}

The most striking finding is that X-VLA's ``entanglement'' disappears entirely under attention-based pruning (99.2\% at $k=0.5$ vs L2-norm's 42.6\%).
This implies the BART encoder does not create irreversible token dependencies; rather, it produces a norm distribution where L2 is a poor importance proxy.
The practical implication is immediate: before concluding that an architecture is ``unprunable,'' practitioners should evaluate multiple importance metrics.

\subsection{The Four-Dimensional Pruning Landscape}

Our data demonstrates that pruning behavior is a function of four interacting factors: (connector, metric, $k$, task).
Attempts to reduce this to fewer dimensions---e.g., predicting behavior from connector type alone---achieve only 62\% accuracy.
The CV scaling law, while valid within Object suite (Spearman $\rho=-1.0$), breaks down cross-suite (mean $|\rho|\approx 0.45$) due to task-specific reversals on Long-horizon.

\subsection{VTR as Architecture-Specific Solution}

VTR's success on SmolVLA (0\%$\rightarrow$59\% at any $M$) and failure on OpenVLA (0\% across all $M$) confirms it specifically repairs pixel-shuffle spatial entanglement rather than providing general compression.
The $M=96$ dip ($p=0.034$) and elevated high-$M$ variance warrant further investigation of training dynamics.

\subsection{Limitations}

\textbf{Simulation only}: All evaluations use LIBERO in simulation; real-robot transfer is untested.
\textbf{MLP class coverage}: Only OpenVLA has full SR evaluation among MLP-projector models (OFT blocked by GPU incompatibility).
\textbf{Missing connector families}: Perceiver Resampler (RoboFlamingo) and Q-Former architectures are not evaluated due to infrastructure constraints.
\textbf{Attention pruning scope}: Attention-based metrics are only fully evaluated on X-VLA; generalization to other architectures requires additional experiments.

\section{Conclusion}

We present the first systematic study of token pruning across VLA architectures, evaluating 5 models $\times$ 4 task suites with 3 pruning methods.
Our findings overturn the assumption that connector architecture alone determines pruning robustness:
pixel-shuffle entanglement is real and irreversible (0\% SR at 5\% removal);
importance metric selection produces larger performance differences (56.6pp) than architectural variation;
and VTR achieves near-lossless compression specifically for pixel-shuffle models while failing on MLP architectures.
These results establish that safe token pruning requires evaluating the full (architecture, metric, ratio, task) configuration space rather than relying on architectural heuristics.

\bibliographystyle{plainnat}
\begin{thebibliography}{10}
\bibitem{openvla} Kim et al. OpenVLA: An Open-Source Vision-Language-Action Model. arXiv:2406.09246, 2024.
\bibitem{pi0} Black et al. $\pi$0: A Vision-Language-Action Flow Model for General Robot Control. arXiv:2410.24164, 2024.
\bibitem{smolvla} Roux et al. SmolVLA: A Vision-Language-Action Model for Affordable Robots. 2025.
\bibitem{xvla} Zhen et al. X-VLA: Cross-Modal Vision-Language-Action Model. 2025.
\bibitem{libero} Liu et al. LIBERO: Benchmarking Knowledge Transfer for Lifelong Robot Learning. NeurIPS 2023.
\bibitem{fastv} Chen et al. An Image is Worth 1/2 Tokens After Layer 2: Plug-and-Play Inference Acceleration for Large Vision-Language Models. ECCV 2024.
\bibitem{efficientvla} EfficientVLA: Training-Free Vision Token Pruning for Efficient VLAs. 2025.
\bibitem{vlaiap} VLA-IAP: Inference-Aware Pruning for Vision-Language-Action Models. 2025.
\bibitem{lightvla} LightVLA: Efficient Vision-Language-Action Models via Light Token Pruning. ICRA 2026.
\bibitem{seewhatmatters} See What Matters: Token Importance Estimation for VLAs. ICML 2026.
\end{thebibliography}

\end{document}
